\documentclass[12pt]{article}

\usepackage[margin=1in]{geometry}
\usepackage{amsmath, amssymb, amsthm}
\usepackage[T1]{fontenc}
\usepackage{lmodern}
\usepackage{microtype}
\usepackage{setspace}
\usepackage{booktabs}
\usepackage{array}
\usepackage{caption}
\usepackage[colorlinks=true, allcolors=blue]{hyperref}

\setlength{\parindent}{0pt}
\setlength{\parskip}{0.6em}
\onehalfspacing

\captionsetup{
    font=small,
    labelfont=bf,
    justification=centering
}

\newtheorem{theorem}{Theorem}[section]
\newtheorem{lemma}[theorem]{Lemma}
\newtheorem{proposition}[theorem]{Proposition}
\newtheorem{corollary}[theorem]{Corollary}

\theoremstyle{definition}
\newtheorem{definition}[theorem]{Definition}

\theoremstyle{remark}
\newtheorem{remark}[theorem]{Remark}

\title{\textbf{Scott's Radiation Activation Law}}
\author{Jaxson Scottlevoy Freeman\\
\small Independent student researcher\\
\small Corresponding author: \texttt{jaxsonstormy@gmail.com}}
\date{October 2026}

\begin{document}

\maketitle

\begin{abstract}
This paper investigates a nonlinear differential equation
with a factorially normalized, Erlang-type forcing term:
\[
\frac{dN}{dt}
=
G_0\frac{\lambda^m t^{m-1}e^{-\lambda t}}{(m-1)!}
-k_1N-\beta N^2.
\]
Here, \(N(t)\) represents a time-dependent state variable,
the first term describes a transient external input, and
the remaining terms represent linear and quadratic losses.
We establish existence, uniqueness, nonnegativity, and
global boundedness of solutions under natural assumptions
on the parameters and initial condition. We derive an
integral representation, obtain an explicit solution in
the absence of quadratic loss, and establish bounds for
the accumulated input and the resulting state variable.
We also examine the long-time behavior of the model and
derive a finite-time approximation framework. The results
provide a mathematical foundation for studying systems
subject to a transient input and nonlinear dissipation.
\end{abstract}

\noindent\textbf{Keywords:} nonlinear differential equation;
Erlang distribution; factorial normalization; existence
and uniqueness; bounded solutions; nonlinear dissipation.

\medskip
\noindent\textbf{2020 Mathematics Subject Classification:}
Primary 34A12; Secondary 34A34, 34C11.

\section{Introduction}

Differential equations provide a general framework for
describing systems whose states change in response to
external inputs and internal loss mechanisms. In many
applications, an external input is concentrated over a
finite period rather than remaining constant indefinitely.
The subsequent evolution may depend on both linear and
nonlinear dissipation.

This paper considers the initial-value problem
\begin{equation}
\frac{dN}{dt}
=
G_0\frac{\lambda^m t^{m-1}e^{-\lambda t}}{(m-1)!}
-k_1N-\beta N^2,
\qquad t\geq 0,
\label{eq:model}
\end{equation}
subject to
\begin{equation}
N(0)=N_0.
\label{eq:initial-condition}
\end{equation}

The forcing term has the form of a scaled Erlang probability
density. The parameters determine the amplitude and timing
of the input, while the terms proportional to \(N\) and
\(N^2\) describe linear and quadratic losses, respectively.

The main questions are whether the initial-value problem
has a unique solution, whether nonnegative initial data
remain nonnegative, whether solutions can grow without
bound, and how the model behaves over long time intervals.
We also investigate the special case in which the quadratic
loss is absent and establish estimates useful for numerical
approximation.

The analysis uses standard methods from ordinary
differential equations, including local existence and
uniqueness, comparison arguments, and integral
representations; see Coddington and Levinson~\cite{coddington}
and Teschl~\cite{teschl}.

\section{Definition and Basic Properties}

\subsection{The mathematical model}

\begin{definition}[Factorial-Erlang Forcing Model]
Let \(m\) be a positive integer and let
\(G_0\geq 0\), \(\lambda>0\), \(k_1\geq 0\),
\(\beta\geq 0\), and \(N_0\geq 0\).
The model is the initial-value problem
\[
\begin{cases}
\displaystyle
\frac{dN}{dt}
=
G_0\frac{\lambda^m t^{m-1}e^{-\lambda t}}{(m-1)!}
-k_1N-\beta N^2,
& t\geq 0,\\[1.2em]
N(0)=N_0.
\end{cases}
\]
\end{definition}

The variables and parameters have the following roles:

\begin{itemize}
    \item \(N(t)\) is the state variable.
    \item \(N_0\) is its initial value.
    \item \(G_0\) controls the total magnitude of the input.
    \item \(\lambda\) is the rate parameter of the input.
    \item \(m\) is a positive integer controlling the shape
    of the forcing profile.
    \item \(k_1\) is the linear loss coefficient.
    \item \(\beta\) is the quadratic loss coefficient.
\end{itemize}

The nonnegative parameter assumptions are appropriate
when \(N\) represents a population, concentration, or
other nonnegative quantity subject to losses. Different
applications may require different parameter assumptions.

For convenience, define the forcing function
\begin{equation}
F_m(t)
=
G_0\frac{\lambda^m t^{m-1}e^{-\lambda t}}{(m-1)!}.
\label{eq:forcing}
\end{equation}
Equation~\eqref{eq:model} can then be written as
\begin{equation}
N'(t)=F_m(t)-k_1N(t)-\beta N(t)^2.
\label{eq:compact-model}
\end{equation}

\subsection{Properties of the forcing term}

The forcing function is nonnegative for \(t\geq 0\).
For \(G_0>0\), it is positive for \(t>0\). Its total
integral is
\begin{equation}
\int_0^\infty F_m(t)\,dt=G_0.
\label{eq:total-input}
\end{equation}

Indeed, the substitution \(u=\lambda t\) gives
\[
\int_0^\infty F_m(t)\,dt
=
\frac{G_0}{(m-1)!}
\int_0^\infty u^{m-1}e^{-u}\,du.
\]
Using the factorial identity
\[
\int_0^\infty u^{m-1}e^{-u}\,du=(m-1)!,
\]
we obtain Equation~\eqref{eq:total-input}.

More generally, the cumulative input delivered by time
\(t\) is
\begin{equation}
A_m(t)
=
\int_0^t F_m(s)\,ds
=
G_0\left[
1-e^{-\lambda t}
\sum_{j=0}^{m-1}\frac{(\lambda t)^j}{j!}
\right].
\label{eq:cumulative-input}
\end{equation}

This identity follows from the standard integral formula
for the Erlang density. In particular,
\[
0\leq A_m(t)\leq G_0,
\qquad
\lim_{t\to\infty}A_m(t)=G_0.
\]

Thus, the forcing term supplies a finite total input,
even though the model is defined over the entire
nonnegative time axis.

\subsection{Interpretation of the loss terms}

The linear loss term is \(-k_1N\). When considered alone,
it produces exponential decay at rate \(k_1\).

The quadratic loss term is \(-\beta N^2\). For
\(\beta>0\), this term increases in magnitude quadratically
with the state variable and can represent density-dependent
loss or other nonlinear dissipation.

When both coefficients vanish, the model reduces to
\(N'=F_m(t)\). When the quadratic coefficient vanishes
but the linear coefficient remains, the equation becomes
linear and can be solved explicitly. The general case
with both loss terms present is nonlinear.

\section{Existence and Uniqueness}

\subsection{Local existence and uniqueness}

\begin{theorem}
For the parameter values specified in the definition,
the initial-value problem
\eqref{eq:model}--\eqref{eq:initial-condition}
has a unique local continuously differentiable solution.
\end{theorem}

\begin{proof}
Define
\[
f(t,N)=F_m(t)-k_1N-\beta N^2.
\]
Since \(m\) is a positive integer and \(\lambda>0\),
the function \(F_m\) is continuous on \([0,\infty)\).
Consequently, \(f\) is continuous in both variables.

Furthermore,
\[
\frac{\partial f}{\partial N}(t,N)
=-k_1-2\beta N
\]
is continuous for all \(t\geq 0\) and \(N\in\mathbb R\).
Hence \(f\) is locally Lipschitz continuous in \(N\).

The standard existence and uniqueness theorem for
ordinary differential equations therefore gives a unique
local solution through the initial point
\((0,N_0)\)~\cite{coddington,teschl}.
\end{proof}

Local existence does not, by itself, guarantee that a
solution exists for every \(t\geq 0\). We establish that
stronger conclusion by proving nonnegativity and
boundedness.

\subsection{Nonnegativity of solutions}

\begin{theorem}
Suppose \(G_0\geq 0\), \(\lambda>0\), \(k_1\geq 0\),
\(\beta\geq 0\), and \(N_0\geq 0\). Every solution of
the initial-value problem remains nonnegative throughout
its interval of existence.
\end{theorem}

\begin{proof}
The equation can be written in the form
\[
N'(t)+\bigl(k_1+\beta N(t)\bigr)N(t)=F_m(t).
\]

Multiplying by the integrating factor associated with
the coefficient \(k_1+\beta N(t)\) gives
\begin{align}
N(t)
={}&N_0
\exp\left(
-\int_0^t[k_1+\beta N(u)]\,du
\right)
\nonumber\\
&+
\int_0^t F_m(s)
\exp\left(
-\int_s^t[k_1+\beta N(u)]\,du
\right)ds.
\label{eq:positive-representation}
\end{align}

This identity holds along any existing solution. The
exponential factors are positive, and \(F_m(s)\geq 0\).
Since \(N_0\geq 0\), both terms on the right-hand side
are nonnegative. Therefore \(N(t)\geq 0\).
\end{proof}

The representation in Equation~\eqref{eq:positive-representation}
is an identity along the solution, rather than an explicit
closed-form solution, because the exponent still depends
on \(N\).

\subsection{Global existence and boundedness}

\begin{theorem}
Under the same parameter assumptions, the unique solution
exists for every \(t\geq 0\) and satisfies
\begin{equation}
0\leq N(t)\leq N_0+A_m(t)\leq N_0+G_0.
\label{eq:global-bound}
\end{equation}
\end{theorem}

\begin{proof}
Nonnegativity has already been established. Therefore,
the loss terms satisfy
\[
-k_1N(t)-\beta N(t)^2\leq 0.
\]
It follows that
\[
N'(t)\leq F_m(t).
\]

Integrating from \(0\) to \(t\) gives
\[
N(t)-N_0\leq\int_0^t F_m(s)\,ds=A_m(t).
\]
Equation~\eqref{eq:cumulative-input} yields
\[
N(t)\leq N_0+A_m(t)\leq N_0+G_0.
\]

Thus the solution remains in the bounded interval
\([0,N_0+G_0]\). The right-hand side of the differential
equation is continuous in time and locally Lipschitz
continuous in the state variable. Since the solution
cannot become unbounded in finite time, the continuation
theorem for ordinary differential equations implies
that it extends to every \(t\geq 0\).
\end{proof}

This theorem establishes global existence without
requiring an explicit formula for the nonlinear solution.

\section{Integral Representation and Special Cases}

\subsection{An integral identity}

Integrating Equation~\eqref{eq:compact-model} from
\(0\) to \(t\) gives the identity
\begin{equation}
\begin{aligned}
N(t)
={}&N_0+A_m(t)
-k_1\int_0^t N(s)\,ds\\
&-\beta\int_0^t N(s)^2\,ds.
\end{aligned}
\label{eq:integral-identity}
\end{equation}

This representation separates the accumulated input
from the cumulative linear and quadratic losses. It also
provides a useful starting point for numerical methods
and analytical estimates.

The identity is exact, but it is not an explicit solution
in the general nonlinear case because the integrals
depend on the unknown state \(N\).

\subsection{The case of no quadratic loss}

When \(\beta=0\), Equation~\eqref{eq:model} becomes
\[
N'(t)+k_1N(t)=F_m(t).
\]

Multiplication by \(e^{k_1t}\) and integration yield
the explicit solution
\begin{equation}
N(t)
=
N_0e^{-k_1t}
+
\int_0^t e^{-k_1(t-s)}F_m(s)\,ds.
\label{eq:linear-solution}
\end{equation}

Substituting the forcing function gives
\begin{align}
N(t)
={}&N_0e^{-k_1t}
\nonumber\\
&+
\frac{G_0\lambda^m e^{-k_1t}}{(m-1)!}
\int_0^t
s^{m-1}e^{-(\lambda-k_1)s}\,ds.
\label{eq:linear-solution-expanded}
\end{align}

This formula remains valid when \(k_1=\lambda\).
In that case, the integral simplifies to
\[
\int_0^t s^{m-1}\,ds=\frac{t^m}{m},
\]
so that
\begin{equation}
N(t)
=
N_0e^{-\lambda t}
+
\frac{G_0\lambda^m t^m e^{-\lambda t}}{m!}.
\label{eq:resonant-linear-solution}
\end{equation}

When \(k_1\neq\lambda\), the integral in
Equation~\eqref{eq:linear-solution-expanded} can also
be expressed using the incomplete gamma function.

\subsection{The case of no losses}

If \(k_1=\beta=0\), then
\[
N'(t)=F_m(t).
\]
Using the initial condition and integrating directly,
we obtain
\begin{equation}
N(t)
=
N_0+
G_0\left[
1-e^{-\lambda t}
\sum_{j=0}^{m-1}\frac{(\lambda t)^j}{j!}
\right].
\label{eq:no-loss-solution}
\end{equation}

In particular,
\[
\lim_{t\to\infty}N(t)=N_0+G_0.
\]
This limiting value represents the initial state plus
the total accumulated input when no loss mechanism
is present.

\subsection{The general nonlinear case}

For \(\beta>0\), the quadratic term makes the equation
nonlinear. The substitution
\[
N(t)=\frac{u'(t)}{\beta u(t)}
\]
transforms the equation, wherever the substitution is
defined, into
\begin{equation}
u''(t)+k_1u'(t)-\beta F_m(t)u(t)=0.
\label{eq:linearized-second-order}
\end{equation}

To verify this transformation, differentiation gives
\[
N'(t)
=
\frac{u''(t)}{\beta u(t)}
-\beta N(t)^2.
\]
Substitution into Equation~\eqref{eq:model} cancels
the quadratic terms and yields
Equation~\eqref{eq:linearized-second-order}.

For \(\beta>0\), the initial conditions corresponding
to \(N(0)=N_0\) can be chosen as
\[
u(0)=1,\qquad u'(0)=\beta N_0.
\]
The resulting solution must be considered on intervals
where \(u(t)\neq 0\).

This transformation provides an alternative linear
second-order formulation, although it does not generally
produce an elementary closed-form expression for \(N(t)\).

\section{Long-Time Behavior}

\subsection{Decay of the forcing term}

For every fixed positive integer \(m\),
\[
\lim_{t\to\infty}F_m(t)=0.
\]
Indeed, exponential decay dominates the polynomial
factor \(t^{m-1}\).

Moreover, the total input is finite:
\[
\int_0^\infty F_m(t)\,dt=G_0.
\]
Thus, the external forcing becomes small at late times,
and its remaining cumulative contribution tends to zero.

\subsection{Decay when linear loss is present}

\begin{theorem}
If \(k_1>0\), then every solution satisfying the
parameter assumptions in Section~2 obeys
\[
\lim_{t\to\infty}N(t)=0.
\]
\end{theorem}

\begin{proof}
Because \(N(t)\geq 0\) and \(\beta\geq 0\),
\[
N'(t)+k_1N(t)
=F_m(t)-\beta N(t)^2
\leq F_m(t).
\]
Multiplying by \(e^{k_1t}\) and integrating yields
\[
N(t)\leq
N_0e^{-k_1t}
+\int_0^t e^{-k_1(t-s)}F_m(s)\,ds.
\]

The first term tends to zero because \(k_1>0\).
For the integral, fix \(T>0\) and split it as
\[
\int_0^t e^{-k_1(t-s)}F_m(s)\,ds
=
\int_0^T e^{-k_1(t-s)}F_m(s)\,ds
+
\int_T^t e^{-k_1(t-s)}F_m(s)\,ds
\]
for \(t\geq T\).

The first integral tends to zero as \(t\to\infty\),
since its interval is fixed. The second is bounded by
\[
\int_T^\infty F_m(s)\,ds,
\]
which can be made arbitrarily small by taking \(T\)
sufficiently large. Therefore the entire integral
tends to zero, proving the result.
\end{proof}

\subsection{Decay when only quadratic loss is present}

\begin{theorem}
Suppose \(k_1=0\), \(\beta>0\), and the other parameter
assumptions remain in force. Then
\[
\lim_{t\to\infty}N(t)=0.
\]
\end{theorem}

\begin{proof}
The equation becomes
\[
N'(t)=F_m(t)-\beta N(t)^2.
\]
Since \(N(t)\geq 0\), integration gives
\[
\beta\int_0^t N(s)^2\,ds
=
N_0+A_m(t)-N(t)
\leq N_0+G_0.
\]
Consequently,
\[
\int_0^\infty N(s)^2\,ds
\leq \frac{N_0+G_0}{\beta}<\infty.
\]

Also,
\[
|N'(t)|\leq F_m(t)+\beta(N_0+G_0)^2,
\]
so \(N'\) is bounded. Hence \(N\) is uniformly
continuous on \([0,\infty)\).

A nonnegative, uniformly continuous function whose
square is integrable on \([0,\infty)\) must tend to
zero. Otherwise, there would be a positive number
\(\varepsilon\) and a sequence of arbitrarily late
times at which \(N(t)\geq\varepsilon\). Uniform
continuity would then force intervals of uniformly
positive length on which \(N\geq\varepsilon/2\),
contradicting square integrability. Therefore
\(N(t)\to 0\).
\end{proof}

Together, these results show that either form of
nonnegative dissipation is sufficient to drive the
state toward zero after the transient input has passed.
If both loss coefficients vanish, the limiting value
instead equals \(N_0+G_0\), as shown in
Equation~\eqref{eq:no-loss-solution}.

\section{Numerical Approximation}

\subsection{A finite-time integral estimate}

Equation~\eqref{eq:integral-identity} gives
\[
N(t)=N_0+A_m(t)
-k_1\int_0^t N(s)\,ds
-\beta\int_0^t N(s)^2\,ds.
\]
Since all terms are nonnegative except for the
explicit subtraction signs, it follows that
\[
0\leq N(t)\leq N_0+A_m(t).
\]

This bound can be used to check numerical calculations.
A computed trajectory that becomes negative or exceeds
\(N_0+A_m(t)\) under the stated parameter assumptions
is inconsistent with the analytical results.

\subsection{An Euler approximation}

For a step size \(h>0\), define \(t_j=jh\) and
approximate the solution using the explicit Euler method:
\begin{equation}
\begin{aligned}
N_{j+1}
={}&N_j+h\Bigg[
G_0\frac{\lambda^m t_j^{m-1}e^{-\lambda t_j}}
{(m-1)!}\\
&\hspace{2.4em}
-k_1N_j-\beta N_j^2
\Bigg].
\end{aligned}
\label{eq:euler}
\end{equation}

The initial approximation is \(N_0\). For sufficiently
small step sizes on a fixed finite interval, the explicit
Euler method converges to the exact solution under the
standard regularity assumptions for this differential
equation.

However, the explicit Euler method does not automatically
preserve nonnegativity for arbitrary step sizes. Numerical
implementations should therefore use an appropriate step
size or a numerical method designed to preserve the
relevant qualitative properties.

\subsection{Local truncation error}

Let
\[
f(t,N)=F_m(t)-k_1N-\beta N^2.
\]
For a sufficiently smooth exact solution, Taylor's
theorem gives
\[
N(t+h)=N(t)+hN'(t)+\frac{h^2}{2}N''(\xi)
\]
for some \(\xi\in(t,t+h)\).

Since \(N'(t)=f(t,N(t))\), the Euler update differs
from the exact solution by a local truncation error
of order \(O(h^2)\), provided the relevant second
derivative is bounded on the interval under
consideration. The resulting global error on a fixed
finite interval is of order \(O(h)\).

These statements concern finite-time numerical
approximation; they do not assert uniform accuracy over
the entire infinite time axis.

\section{Possible Generalizations}

The forcing term can be modified by changing its amplitude,
rate parameter, or shape. For example, one may consider
\begin{equation}
N'(t)
=
G_0\frac{\lambda^m t^{m-1}e^{-\lambda t}}{(m-1)!}
-k_1N(t)-\beta N(t)^p,
\label{eq:generalized-model}
\end{equation}
where \(p>1\).

For nonnegative solutions and nonnegative loss
coefficients, the same basic upper bound
\[
N(t)\leq N_0+A_m(t)\leq N_0+G_0
\]
continues to hold, provided the generalized problem
has the required existence and uniqueness properties.
For real \(p\) that is not an integer, the expression
\(N^p\) is naturally defined here on \(N\geq 0\).

Another variation replaces the single Erlang-type input
with a finite sum of such inputs:
\[
F(t)=
\sum_{r=1}^{q}
G_r\frac{\lambda_r^{m_r}t^{m_r-1}
e^{-\lambda_r t}}{(m_r-1)!},
\]
where \(q\) and the \(m_r\) are positive integers,
\(G_r\geq 0\), and \(\lambda_r>0\).

In this setting,
\[
\int_0^\infty F(t)\,dt=\sum_{r=1}^{q}G_r.
\]
The boundedness argument then gives
\[
0\leq N(t)\leq N_0+\sum_{r=1}^{q}G_r
\]
for nonnegative loss coefficients and nonnegative
initial data.

These extensions allow the model to represent several
distinct transient inputs or different forms of
nonlinear dissipation. Their practical usefulness
depends on the system being modeled and on the
interpretation of the parameters.

\section{Conclusion}

This paper examines a nonlinear differential equation
with a factorially normalized, Erlang-type forcing term
and linear and quadratic loss mechanisms. Under natural
nonnegativity assumptions, the initial-value problem
has a unique global solution. The solution remains
nonnegative and satisfies the bound
\[
0\leq N(t)\leq N_0+G_0.
\]

The cumulative input has an explicit expression, and
the model admits an explicit solution when quadratic
loss is absent. In the general nonlinear case, an
integral representation and a transformation to a
second-order linear equation provide alternative ways
to analyze the dynamics.

The long-time results show that the state tends to zero
when either linear or quadratic dissipation is present.
In the absence of both loss mechanisms, the state
approaches the initial value plus the total input.
The analytical bounds also provide consistency checks
for numerical approximations.

Further work could investigate parameter estimation,
sensitivity to the shape parameter \(m\), improved
numerical schemes, and applications to specific physical
or biological systems. Such investigations would help
determine when this model is appropriate and which
features of its solutions are most relevant in practice.

\section{Acknowledgments}

The author thanks the teachers and mentors who encouraged
the pursuit of mathematical exploration.

\begin{thebibliography}{9}

\bibitem{coddington}
E. A. Coddington and N. Levinson,
\emph{Theory of Ordinary Differential Equations},
McGraw-Hill, New York, 1955.

\bibitem{teschl}
G. Teschl,
\emph{Ordinary Differential Equations and Dynamical Systems},
Graduate Studies in Mathematics, vol. 140,
American Mathematical Society, Providence, RI, 2012.

\end{thebibliography}

\end{document}